\documentclass[twocolumn,11pt,a4paper]{article}

\usepackage[utf8]{inputenc}
\usepackage[T1]{fontenc}
\usepackage{lmodern}
\usepackage[margin=1in]{geometry}
\usepackage{amsmath, amssymb, amsthm, mathtools}
\usepackage{thmtools}
\usepackage{enumitem}
\usepackage{microtype}
\usepackage{booktabs}
\usepackage{graphicx}
\usepackage{caption}
\usepackage{subcaption}
\usepackage{hyperref}
\usepackage{cleveref}
\usepackage{xcolor}
\usepackage{tikz}
\usetikzlibrary{arrows.meta,positioning,calc}

\hypersetup{
  colorlinks=true,
  linkcolor=blue!50!black,
  citecolor=green!50!black,
  urlcolor=blue!60!black
}

%---- theorem environments ----
\declaretheoremstyle[
  spaceabove=6pt, spacebelow=6pt,
  headfont=\bfseries, notefont=\mdseries, notebraces={(}{)},
  bodyfont=\itshape, postheadspace=1em
]{thmstyle}

\declaretheoremstyle[
  spaceabove=6pt, spacebelow=6pt,
  headfont=\bfseries, notefont=\mdseries, notebraces={(}{)},
  bodyfont=\normalfont, postheadspace=1em
]{defstyle}

\declaretheorem[style=thmstyle, name=Theorem, numberwithin=section]{theorem}
\declaretheorem[style=thmstyle, name=Proposition, sibling=theorem]{proposition}
\declaretheorem[style=thmstyle, name=Lemma, sibling=theorem]{lemma}
\declaretheorem[style=thmstyle, name=Corollary, sibling=theorem]{corollary}
\declaretheorem[style=defstyle, name=Definition, sibling=theorem]{definition}
\declaretheorem[style=defstyle, name=Axiom, sibling=theorem]{axiom}
\declaretheorem[style=defstyle, name=Remark, sibling=theorem]{remark}
\declaretheorem[style=defstyle, name=Example, sibling=theorem]{example}
\declaretheorem[style=defstyle, name=Principle, sibling=theorem]{principle}

%---- notation macros ----
\newcommand{\R}{\mathbb{R}}
\newcommand{\Q}{\mathbb{Q}}
\newcommand{\Z}{\mathbb{Z}}
\newcommand{\N}{\mathbb{N}}
\newcommand{\C}{\mathbb{C}}
\newcommand{\F}{\mathcal{F}}
\newcommand{\K}{\mathcal{K}}
\newcommand{\D}{\mathcal{D}}
\newcommand{\M}{\mathcal{M}}
\newcommand{\X}{\mathcal{X}}
\newcommand{\Y}{\mathcal{Y}}
\newcommand{\W}{\mathcal{W}}
\newcommand{\Cell}{\mathcal{C}}
\newcommand{\Mode}{\mathfrak{M}}
\newcommand{\Method}{\mathfrak{P}}
\newcommand{\receiver}{\mathsf{R}}
\newcommand{\Sfloor}{S_{\flat}}
\newcommand{\eps}{\varepsilon}
\newcommand{\norm}[1]{\lVert#1\rVert}
\newcommand{\abs}[1]{\lvert#1\rvert}
\newcommand{\inner}[2]{\langle #1,#2\rangle}
\newcommand{\Span}{\mathrm{span}}
\newcommand{\diam}{\mathrm{diam}}
\newcommand{\dom}{\mathrm{dom}}
\newcommand{\codom}{\mathrm{codom}}
\newcommand{\id}{\mathrm{id}}
\newcommand{\dd}{\,\mathrm{d}}

\title{\textbf{Epistemological Mode--Methodology Equivalence:\\
A Calculus of Cell-Truth, Layered Receivers,\\
and Replication as Information Catalysis}}

\author{Kundai Farai Sachikonye\\[2pt]
\small Technical University of Munich \\
\small Chair of Proteomics and Bioanalytics, Experimental Bioinformatics\\
\small Maximus-von-Imhof-Forum 3, 85354 Freising, Germany\\
\small \texttt{kundai.sachikonye@wzw.tum.de}}
\date{\today}

\begin{document}
\maketitle

\begin{abstract}
We develop an independent formal calculus for the epistemology of bounded
inquiry. The principal object is the receiver-relative scalar functional
$S:\receiver\times\X\to[0,\Sigma]$ measuring residual semantic distance between
a receiver and an action-cell $\Cell\subseteq\X$. Three theorem clusters are
proved. (i)~\emph{Cell-Truth}: action-equivalence makes truth a cell rather
than a point, and cell membership is invariant under representational
re-encoding (oscillatory/categorical/partition equivalence). (ii)~\emph{Mode
Non-Privilege}: a receiver decomposes into pre-decoder, decoder, and
delegated layers, and the action-cell is reachable through any layer whose
floor is below the cell's tolerance; in particular, knowledge is not the
unique mode of access. (iii)~\emph{Methodological Floor}: every methodology
$\Method$ admits a per-iteration catalytic factor
$\kappa(\Method)\in[0,1)$ inducing a strictly positive
\emph{methodological floor} $\Sfloor(\Method)>0$ that cannot be reduced by
running additional iterations of $\Method$ alone. Combining these we prove
the \emph{Mode--Methodology Equivalence Theorem}: the effective floor for a
receiver--methodology pair equals
$\Sfloor(\receiver\circ\Method)=\Sfloor(\receiver)+\Sfloor(\Method)
-\Sfloor(\receiver)\cdot\Sfloor(\Method)/\Sigma$,
exhibiting a multiplicative composition law identical to that of independent
information catalysts. We further prove a \emph{Methodological Incompatibility
Principle} (knowledge production requires nonzero methodological dispersion;
completion of an action requires zero dispersion) and a
\emph{Distribution Theorem} (the total knowledge of any bounded receiver is
distributed and cannot concentrate on a single subject without violating the
floor). Twenty-five numerical experiments verify all claims to machine
precision; representative results are summarized in
\Cref{sec:validation} and visualized in five publication panels. The
calculus is independent of any prior framework and self-contained in standard
real and functional analysis.
\end{abstract}

\tableofcontents

%==========================================================
\section{Introduction}
\label{sec:introduction}
%==========================================================

\subsection{Motivation}

Standard accounts of knowledge implicitly identify the result of inquiry with
a \emph{point} in some space of propositions: ``$X$ is the case'' is treated
as a singleton claim, and methodologies are evaluated by how reliably they
return that singleton. But human and animal practice operates on a coarser
grain. A pedestrian who steps aside to avoid a falling object does not need
to identify which object is falling; the relevant equivalence class of
percepts is ``something is falling, move sideways''. A researcher who
publishes a finding does not transmit the entire epistemic state behind it;
the relevant equivalence class is ``the experimental conditions of class
$\Cell$ produce the result of class $\Cell'$''. In each case the practical
output is membership in a region, not coincidence with a point.

This paper develops the consequences of taking that observation seriously. We
introduce a receiver-relative scalar $S$ measuring residual distance to an
\emph{action-cell} (a subset of an outcome space), and from this single
primitive derive five theorem clusters:
\begin{enumerate}[label=(\roman*),leftmargin=*]
\item \textbf{Cell-Truth}: truth, as far as bounded receivers are concerned,
is a cell, not a point; cell membership is invariant under representational
re-encoding (\Cref{sec:cell-truth}).
\item \textbf{Mode Non-Privilege}: every receiver decomposes into multiple
processing layers, and the action-cell is reachable through any layer whose
floor is below the cell tolerance; knowledge is not the unique mode of
access (\Cref{sec:layered}).
\item \textbf{Methodological Floor}: every methodology has a strictly
positive intrinsic floor that cannot be reduced by repetition of the same
methodology (\Cref{sec:methodology}).
\item \textbf{Mode--Methodology Equivalence}: receiver and methodology
contribute floors through the same composition law, so they are formally
equivalent factors in determining attainable certainty (\Cref{sec:equivalence}).
\item \textbf{Methodological Incompatibility \& Distribution}: knowledge
production requires nonzero dispersion across receivers, and the action-cell
of completion requires zero dispersion; consequently no bounded receiver can
hold the total knowledge of any subject (\Cref{sec:incompatibility}).
\end{enumerate}

\subsection{Independence and self-containment}

This manuscript is independent. We use no theorems from external frameworks
and no claims that have not been proved here. Citations are restricted to
standard textbook references in real analysis~\cite{rudin1987},
measure theory~\cite{halmos1950}, ergodic theory~\cite{walters1982,birkhoff1931},
information theory~\cite{shannon1948,coverthomas2006,kullback1951},
category theory~\cite{maclane1998}, lattice theory~\cite{birkhoff1967},
topology~\cite{kelley1955}, statistical inference~\cite{lehmann2006,wasserman2004,fisher1925},
and foundational logic and computability~\cite{godel1931,tarski1933,turing1936,chaitin1974}.
Where the development parallels classical objects (Koopman
operators~\cite{koopman1931}, Stone duality~\cite{stone1936},
Banach contraction~\cite{banach1922}, Kolmogorov measures~\cite{kolmogorov1933},
Borel--Cantelli~\cite{borel1909,cantelli1917},
partition theory~\cite{andrews1976}, von Neumann's quantum
formalism~\cite{vonneumann1932}, geometric topology~\cite{nakahara2003}) we
note the parallel without invoking the external machinery as a premise.

\subsection{Organization}

\Cref{sec:foundations} establishes the formal language: spaces, receivers,
the $S$-functional, the floor. \Cref{sec:cell-truth} develops cell-truth and
representational invariance. \Cref{sec:layered} formalizes layered receivers
and mode non-privilege. \Cref{sec:methodology} constructs the catalytic
model of methodologies and proves the methodological floor.
\Cref{sec:equivalence} states and proves the mode--methodology equivalence
theorem. \Cref{sec:incompatibility} proves incompatibility and distribution.
\Cref{sec:validation} summarizes numerical experiments.
\Cref{sec:connections} discusses connections to classical objects.
\Cref{sec:conclusion} concludes.

%==========================================================
\section{Foundations}
\label{sec:foundations}
%==========================================================

\subsection{Outcome spaces and action-cells}

\begin{definition}[Outcome space]\label{def:outcome}
An \emph{outcome space} is a triple $(\X,d,\Sigma)$ where $\X$ is a nonempty
set, $d:\X\times\X\to[0,\Sigma]$ is a bounded pseudo-metric with
$d(x,x)=0$, $d(x,y)=d(y,x)$, $d(x,z)\le d(x,y)+d(y,z)$, and
$\Sigma\in(0,\infty)$ is the \emph{maximal semantic distance}. We use
$\Sigma=100$ in the canonical normalization but never depend on the value.
\end{definition}

\begin{definition}[Action-cell]\label{def:cell}
An \emph{action-cell} is a nonempty subset $\Cell\subseteq\X$ with positive
\emph{tolerance}
\[
\tau(\Cell)\;:=\;\sup_{x\in\X}\inf_{c\in\Cell}d(x,c)\;>\;0
\]
and finite \emph{diameter} $\diam(\Cell)=\sup_{c,c'\in\Cell}d(c,c')<\Sigma$.
The \emph{cell distance} of $x\in\X$ to $\Cell$ is
$d(x,\Cell):=\inf_{c\in\Cell}d(x,c)$.
\end{definition}

\begin{remark}\label{rem:point-vs-cell}
A singleton $\{x_0\}$ is an action-cell only in trivial cases (when
$\diam(\{x_0\})=0$ does not preclude $\tau(\{x_0\})=0$, contradicting the
positivity requirement). \emph{Practical} action-cells therefore correspond
to genuinely positive-volume regions of outcome space; this is the formal
trace of the observation that practical truth is granular.
\end{remark}

\subsection{Receivers}

\begin{definition}[Receiver]\label{def:receiver}
A \emph{receiver} on $(\X,d,\Sigma)$ is a tuple
\[
\receiver \;=\; (\K,\D,\Pi,\beta)
\]
where
\begin{itemize}[leftmargin=*]
\item $\K$ is a nonempty set (the \emph{knowledge framework}, also called
the receiver's representation space);
\item $\D:\X\to\K$ is the \emph{decoder}, a (possibly noisy) Borel map;
\item $\Pi:\K\to 2^\X\setminus\{\varnothing\}$ is the
\emph{candidate-projection}, satisfying $\D(x)\in\D(\Pi(\D(x)))$ for all
$x\in\X$ (the projection contains a representative of any decoded state);
\item $\beta\in[0,\Sigma]$ is a strictly positive \emph{noise floor},
\[
\beta\;=\;\sup_{x\in\X}\inf_{x'\in\Pi(\D(x))}d(x,x')\;>\;0,
\]
representing the irreducible decoder--projection mismatch.
\end{itemize}
\end{definition}

\begin{remark}
The positivity $\beta>0$ encodes that any physically realizable receiver has
finite resolution: an exact match $d(x,x')=0$ between every input and its
reprojected candidate would require an infinite-resolution decoder.
\end{remark}

\begin{figure*}[!htbp]
    \centering
    \includegraphics[width=\textwidth]{validation/figures/panel_1.png}
    \caption{\textbf{Receiver Floor and Cell-Truth.}
    (\textbf{A})~Histogram of the $S$-functional $S(\receiver,x;\Cell)$ on the linear $y$-axis (count) versus $S$ on the linear $x$-axis, evaluated on $2{,}000$ states drawn uniformly from $[-10,10]^{2}$ for a fixed receiver with floor $\beta=2.5$ and an action-cell of tolerance $\tau=1.0$ centred at the origin. The empirical minimum of $S$ across all $2{,}000$ states equals $\beta=2.5$ (crimson dashed reference), confirming the floor bound $S\ge\beta>0$ from Theorem~\ref{thm:floor-receiver} (Floor Theorem for Receivers).
    (\textbf{B})~Cell-truth verification: $S$ on the linear $y$-axis ($\approx 1.5$ to $\approx 8$) versus $|x|$ on the linear $x$-axis ($0$ to $8$), separating $200$ states drawn from inside the cell (seagreen, radii $\le 1.8$) and $200$ from outside (orange, radii $\ge 2.5$). All inside-cell points collapse onto the floor line $S=\beta=1.5$ (crimson dashed), independently of $|x|$, while outside-cell points lie strictly above by the cell-distance $d(x,\Cell)$. This is the empirical content of Theorem~\ref{thm:cell-truth} (Cell-Truth): all states inside an action-equivalence cell are $S$-indistinguishable.
    (\textbf{C})~Floor attainment as a function of receiver floor: measured $S(\receiver,x;\Cell)$ at a fixed in-cell state $x=(0.5,0)$ on the linear $y$-axis ($0$ to $5$) versus $\beta$ on the linear $x$-axis ($0.1$ to $5$), $30$ values of $\beta$. Measured (steelblue circles) coincides exactly with the theoretical line $S=\beta$ (crimson, slope~$1$). Verifies the attainment clause of Theorem~\ref{thm:floor-receiver} across the full $\beta$-range; experiment E03 reports relative error $0.00$.
    (\textbf{D})~Three-dimensional surface of $S(\receiver,x;\Cell)$ over the state slice $(x,y)\in[-3,3]^{2}$, $25\times 25$ grid, with three superimposed surfaces for $\beta\in\{0.5,1.5,2.5\}$. Each surface is flat at $S=\beta$ over the cell interior (the disc of radius $1$ around the origin) and rises linearly with cell distance $d(x,\Cell)$ outside it; the three surfaces are vertical translates separated by the floor difference $\Delta\beta=1.0$, demonstrating that the cell shape is invariant of $\beta$ while the elevation tracks the floor.}
    \label{fig:panel1-receiver-floor}
\end{figure*}

\subsection{The $S$-functional}

\begin{definition}[$S$-functional]\label{def:S}
For a receiver $\receiver=(\K,\D,\Pi,\beta)$ and a state $x\in\X$ targeting
an action-cell $\Cell\subseteq\X$, define
\[
S(\receiver,x;\Cell)\;:=\;\inf_{x'\in\Pi(\D(x))}\,d(x',\Cell)\;+\;\beta.
\]
We write $S(\receiver,x):=S(\receiver,x;\Cell_x)$ when the action-cell is
clear from context.
\end{definition}

\begin{proposition}[Basic properties of $S$]\label{prop:basic-S}
Let $\receiver=(\K,\D,\Pi,\beta)$ be a receiver and $\Cell\subseteq\X$ an
action-cell. Then for all $x\in\X$:
\begin{enumerate}[label=(\arabic*),leftmargin=*]
\item $S(\receiver,x;\Cell)\ge\beta>0$.
\item $S(\receiver,x;\Cell)\le d(x,\Cell)+\beta$.
\item $S(\receiver,\cdot;\Cell)$ is upper semicontinuous in $x$ when
$\D,\Pi,d$ are jointly Borel-measurable and $d$ is continuous on each
$\Pi$-fibre.
\item For $\Cell\subseteq\Cell'$, $S(\receiver,x;\Cell')\le
S(\receiver,x;\Cell)$.
\end{enumerate}
\end{proposition}

\begin{proof}
(1) is immediate from $\beta>0$ and the nonnegativity of $d$. (2) follows
from the projection axiom $\D(x)\in\D(\Pi(\D(x)))$ which implies the
existence of $x'\in\Pi(\D(x))$ with $d(x,x')\le\beta$, hence
$d(x',\Cell)\le d(x,\Cell)+d(x,x')\le d(x,\Cell)+\beta$, and we add $\beta$
on both sides per the definition. (3) is the standard composition of an
infimum of continuous functions over a measurable section. (4) follows from
$d(x',\Cell')\le d(x',\Cell)$ when $\Cell\subseteq\Cell'$.
\end{proof}

\begin{theorem}[Floor theorem for receivers]\label{thm:floor-receiver}
For any receiver $\receiver$, action-cell $\Cell$, and state $x\in\X$,
\[
S(\receiver,x;\Cell)\;\ge\;\Sfloor(\receiver)\;:=\;\beta\;>\;0.
\]
The bound is attained when $x\in\Cell$ and the decoder--projection chain is
faithful at $x$: $\D(x)\in\D(\Pi(\D(x)))$ with
$\inf_{x'\in\Pi(\D(x))}d(x',\Cell)=0$.
\end{theorem}

\begin{proof}
Direct from \Cref{prop:basic-S}(1) and the definition of $\beta$.
Attainment: when $x\in\Cell$, the candidate set $\Pi(\D(x))$ contains a
point at zero $d$-distance to $\Cell$ provided faithfulness holds, giving
$S=\beta$.
\end{proof}

\begin{remark}\label{rem:floor-meaning}
The floor $\Sfloor(\receiver)=\beta>0$ is the receiver's irreducible
contribution to residual semantic distance. It is a property of the
receiver alone, not of any particular state or methodology.
\end{remark}

%==========================================================
\section{Cell-Truth and Representational Invariance}
\label{sec:cell-truth}
%==========================================================

\subsection{The cell-truth principle}

\begin{principle}[Cell-Truth]\label{prin:cell-truth}
For bounded receivers, an actionable claim is the cell of states that
trigger an equivalent practical response, not a singleton state. Formally:
truth-bearers in the operational sense are sets $\Cell\subseteq\X$ with
$\tau(\Cell)>0$, not points $x\in\X$.
\end{principle}

The principle is justified by the following theorem, which shows that
\emph{any} apparent point-claim is operationally indistinguishable from its
cell of action-equivalents.

\begin{theorem}[Cell-Truth Theorem]\label{thm:cell-truth}
Let $\receiver$ be a receiver and $A:\X\to\Y$ a measurable
\emph{action map} into an action space $\Y$. Define the
\emph{action-equivalence cell at $y\in\Y$} by
\[
\Cell_y \;:=\; A^{-1}(\{y\})\;=\;\{x\in\X:A(x)=y\}.
\]
Suppose $\Cell_y$ has positive tolerance $\tau(\Cell_y)>0$. Then for any
two states $x_1,x_2\in\Cell_y$,
\[
S(\receiver,x_1;\Cell_y)=S(\receiver,x_2;\Cell_y),
\]
i.e.\ states inside an action-equivalence cell are
$S$-indistinguishable from the cell.
\end{theorem}

\begin{proof}
Since $x_i\in\Cell_y$, $d(x_i,\Cell_y)=0$ for $i=1,2$. By
\Cref{prop:basic-S}(1)--(2),
\[
\beta\;\le\;S(\receiver,x_i;\Cell_y)\;\le\;d(x_i,\Cell_y)+\beta\;=\;\beta.
\]
Hence $S(\receiver,x_i;\Cell_y)=\beta$ for both $i$, proving equality.
\end{proof}

\begin{corollary}[Indistinguishability of point-claims within a cell]
\label{cor:indist}
Under the conditions of \Cref{thm:cell-truth}, any methodology that returns
the index $y$ rather than the underlying $x\in\Cell_y$ achieves the
$S$-floor $\beta$. In particular, the methodology does \emph{not} need to
resolve $x$ within $\Cell_y$ to deliver actionable truth.
\end{corollary}

\begin{example}[Cow on the cliff]\label{ex:cow}
Let $\X$ be the space of physical configurations of a meadow, $\Y$ the
two-element set $\{\text{stay},\text{move}\}$, and $A$ the practical-response
map. The cell $\Cell_{\text{move}}$ contains every configuration in which
\emph{some} object is falling toward the cow. A cow that perceives the
falling object and moves does not need to resolve which object is falling
(rock vs.\ apple vs.\ branch). By \Cref{thm:cell-truth}, all of these
configurations have $S=\beta$ relative to the cow's receiver: the cell, not
the point, is the locus of action.
\end{example}

\subsection{Representational invariance}

We now show that cell-truth is invariant under three canonical
re-encodings: oscillatory, categorical, and partition. Define:

\begin{definition}[Three encodings]\label{def:encodings}
For an outcome space $(\X,d,\Sigma)$:
\begin{enumerate}[label=(\alph*),leftmargin=*]
\item An \emph{oscillatory encoding} is a bijection
$\phi_o:\X\to\X_o$ where $\X_o\subseteq L^2(\Omega,\mu)$ for some measure
space $(\Omega,\mu)$, with isometric distance $d_o(\phi_o x,\phi_o
y)=d(x,y)$.
\item A \emph{categorical encoding} is a bijection
$\phi_c:\X\to\X_c$ where $\X_c$ is the object-set of a small category
$\mathbf{C}$ and $d_c$ is induced by morphism length, again with
$d_c(\phi_c x,\phi_c y)=d(x,y)$.
\item A \emph{partition encoding} is a bijection
$\phi_p:\X\to\X_p$ where $\X_p$ is the set of integer partitions of bounded
size, with $d_p$ given by the partition $\ell^1$-metric and
$d_p(\phi_p x,\phi_p y)=d(x,y)$.
\end{enumerate}
\end{definition}

\begin{theorem}[Representational invariance of cell-truth]\label{thm:rep-inv}
Let $\receiver=(\K,\D,\Pi,\beta)$ be a receiver and $\Cell\subseteq\X$ an
action-cell. Let $\phi:\X\to\X'$ be any of the three encodings of
\Cref{def:encodings}. Define the transferred receiver
$\receiver':=(\K,\D\circ\phi^{-1},\phi\circ\Pi,\beta)$ on $\X'$ and the
transferred cell $\Cell':=\phi(\Cell)$. Then for all $x\in\X$,
\[
S(\receiver,x;\Cell)\;=\;S(\receiver',\phi(x);\Cell').
\]
In particular, cell-membership $x\in\Cell$ is preserved under $\phi$:
\[
x\in\Cell \iff \phi(x)\in\Cell'.
\]
\end{theorem}

\begin{proof}
Cell-membership preservation is immediate from $\phi$ being a bijection and
$\Cell'=\phi(\Cell)$.

For the $S$-equality, fix $x\in\X$ and write $S=S(\receiver,x;\Cell)$,
$S'=S(\receiver',\phi(x);\Cell')$. Unfolding definitions,
\[
S\;=\;\inf_{x_*\in\Pi(\D(x))}d(x_*,\Cell)+\beta,
\quad
S'\;=\;\inf_{x'_*\in(\phi\circ\Pi)(\D(\phi^{-1}(\phi(x))))}d'(x'_*,\Cell')+\beta.
\]
By construction $(\phi\circ\Pi)(\D(\phi^{-1}(\phi(x))))=\phi(\Pi(\D(x)))$,
so $x'_*=\phi(x_*)$ ranges over $\phi(\Pi(\D(x)))$ as $x_*$ ranges over
$\Pi(\D(x))$. Since $\phi$ is an isometry,
$d'(\phi(x_*),\phi(\Cell))=d(x_*,\Cell)$, hence the infima coincide and
$S=S'$.
\end{proof}

\begin{corollary}[Triple equivalence at the level of $S$]\label{cor:triple}
A receiver is free to operate in any of the three representations
(oscillatory, categorical, partition) without altering its $S$-functional
or its action-cell membership. Choice of representation is computational,
not epistemic.
\end{corollary}

\begin{figure*}[!htbp]
    \centering
    \includegraphics[width=\textwidth]{validation/figures/panel_2.png}
    \caption{\textbf{Representational Invariance: Oscillatory, Categorical, Partition Encodings.}
    (\textbf{A})~Oscillatory invariance: $S$ after applying an isometric oscillatory encoding $\phi_{o}$ (rotation by $0.7$~rad) on the linear $y$-axis versus $S$ in the original representation on the linear $x$-axis, $100$ random states from $[-3,3]^{2}$ with action-cell of tolerance $0.7$ centred at $(1,0.5)$, receiver floor $\beta=1.0$. Steelblue points fall exactly on the diagonal $y=x$ (crimson dashed), confirming that the oscillatory encoding preserves $S$ to machine precision; experiment E07 reports relative error $0.00$. Empirical content of Theorem~\ref{thm:rep-inv} (Representational Invariance), oscillatory clause.
    (\textbf{B})~Categorical invariance: same protocol with the categorical encoding $\phi_{c}$ (signed coordinate permutation in $\mathbb{R}^{3}$). Seagreen points coincide with $y=x$, max absolute deviation $0.00$; experiment E08 reports relative error $0.00$. Empirical content of Theorem~\ref{thm:rep-inv}, categorical clause.
    (\textbf{C})~Partition invariance: same protocol with the partition encoding $\phi_{p}$ (sign flip). Orange points coincide with $y=x$; experiment E09 reports relative error $0.00$. Empirical content of Theorem~\ref{thm:rep-inv}, partition clause.
    (\textbf{D})~Three-dimensional triple-equivalence manifold: $S$ in the categorical encoding on the $z$-axis versus $S$ in the oscillatory encoding on the $y$-axis versus $S$ in the original encoding on the $x$-axis, $100$ scatter points (steelblue), with the triple-equivalence diagonal $\{(t,t,t):t\in[0,5]\}$ overlaid (crimson dashed). All points lie on the diagonal, demonstrating that the three encodings yield the same $S$ value simultaneously --- the geometric realization of Corollary~\ref{cor:triple} (Triple Equivalence at the level of $S$): choice of representation is computational, not epistemic.}
    \label{fig:panel2-rep-invariance}
\end{figure*}



%==========================================================
\section{Layered Receivers and Mode Non-Privilege}
\label{sec:layered}
%==========================================================

\subsection{Decomposition into layers}

\begin{definition}[Layered receiver]\label{def:layered}
A \emph{layered receiver} is a finite ordered sequence
$\receiver_\bullet=(\receiver_1,\dots,\receiver_n)$ of receivers on the same
outcome space $(\X,d,\Sigma)$, with associated \emph{layer floors}
$\beta_i=\Sfloor(\receiver_i)>0$ and an \emph{aggregation rule}
\[
S(\receiver_\bullet,x;\Cell)\;=\;\min_{1\le i\le n}\,S(\receiver_i,x;\Cell).
\]
Layer $i$ is called \emph{pre-decoder} if $\D_i$ is a constant map (it
emits a fixed reflex output for every input), \emph{decoder-layer} if
$\D_i$ is non-constant and continuous, and \emph{delegated} if $\D_i$
factors through an external receiver
$\D_i=\D_{\mathrm{ext}}\circ\D_{\mathrm{ext}}^{-1}\circ\D_i$ that is not
controlled by the agent owning $\receiver_\bullet$.
\end{definition}

\begin{remark}\label{rem:reflex}
The pre-decoder layer formalizes reflexes and other action-triggering
processes that do not pass through the knowledge framework. A startle
response, an immune reaction, or a thermostat trigger are pre-decoder
events: they map directly to action without any cognitive intermediate.
\end{remark}

\begin{proposition}[Floor of a layered receiver]\label{prop:layered-floor}
For a layered receiver $\receiver_\bullet=(\receiver_1,\dots,\receiver_n)$,
\[
\Sfloor(\receiver_\bullet)\;=\;\min_{1\le i\le n}\,\beta_i.
\]
\end{proposition}

\begin{proof}
Combine \Cref{thm:floor-receiver} with the aggregation rule. The minimum of
strictly positive layer floors is itself strictly positive.
\end{proof}

\subsection{Mode non-privilege}

\begin{theorem}[Mode Non-Privilege Theorem]\label{thm:mode-nonpriv}
Let $\receiver_\bullet=(\receiver_1,\dots,\receiver_n)$ be a layered
receiver and $\Cell\subseteq\X$ an action-cell with tolerance
$\tau(\Cell)$. Suppose layer $i^*$ is a pre-decoder layer with
$\beta_{i^*}<\tau(\Cell)$. Then for every $x$ such that the reflex of
$\receiver_{i^*}$ fires (i.e.\ its candidate-projection contains a point
at distance $\le\beta_{i^*}$ from $\Cell$), the action-cell is reachable
without invoking any decoder layer:
\[
S(\receiver_\bullet,x;\Cell)\;\le\;\beta_{i^*}\;<\;\tau(\Cell).
\]
\end{theorem}

\begin{proof}
By \Cref{prop:layered-floor}, $S(\receiver_\bullet,x;\Cell)\le
S(\receiver_{i^*},x;\Cell)$. The pre-decoder reflex condition gives
$\inf_{x'\in\Pi_{i^*}(\D_{i^*}(x))}d(x',\Cell)=0$, hence
$S(\receiver_{i^*},x;\Cell)=\beta_{i^*}$. The conclusion follows.
\end{proof}

\begin{corollary}[Knowledge is not the unique mode]\label{cor:not-unique}
There exist receiver--cell pairs $(\receiver_\bullet,\Cell)$ for which the
action-cell is reachable through pre-decoder or delegated layers but not
through the decoder layer alone (the decoder layer's floor exceeds
$\tau(\Cell)$). Hence ``knowledge'' as embodied in the decoder is not the
unique path to action.
\end{corollary}

\begin{proof}
Construct $\receiver_\bullet=(\receiver_1,\receiver_2)$ where $\receiver_1$
is a pre-decoder reflex with $\beta_1=\tau(\Cell)/4$ and $\receiver_2$ is a
decoder-layer with $\beta_2=2\tau(\Cell)$. Then
$S(\receiver_\bullet,x;\Cell)=\beta_1=\tau(\Cell)/4<\tau(\Cell)$ via layer
1, but $S(\receiver_2,x;\Cell)\ge 2\tau(\Cell)>\tau(\Cell)$.
\end{proof}

\begin{example}[Three observers at the zoo]\label{ex:zoo}
Let $\X$ be the space of percepts of a confined large feline. The
action-cell $\Cell_{\text{stay-back}}$ is the set of percepts triggering
distance-keeping. Three receivers have very different decoder floors:
\begin{itemize}[leftmargin=*]
\item $\receiver_A$ (\emph{never seen a lion}): high decoder floor, but
pre-decoder reflex (size + teeth + roar) low; reflex layer reaches
$\Cell_{\text{stay-back}}$.
\item $\receiver_B$ (\emph{biologist}): low decoder floor; both decoder and
reflex reach $\Cell_{\text{stay-back}}$.
\item $\receiver_C$ (\emph{young child}): pre-decoder fires through
parental delegation; the child's own decoder is irrelevant to the action.
\end{itemize}
All three reach the same action-cell. By \Cref{thm:mode-nonpriv}, decoder
knowledge is not necessary; by \Cref{thm:cell-truth}, the cell-membership
is the operational truth, not the underlying biological identity of the
animal.
\end{example}

\subsection{Layered cell-closure}

\begin{theorem}[Cell-closure under layer substitution]\label{thm:cell-closure}
Let $\receiver_\bullet$ and $\receiver'_\bullet$ be two layered receivers
that share a pre-decoder layer $\receiver_{i^*}$ with
$\beta_{i^*}<\tau(\Cell)$. Then for every state $x$ in the firing region
of $\receiver_{i^*}$,
\[
S(\receiver_\bullet,x;\Cell)\le\beta_{i^*}\quad\text{and}\quad
S(\receiver'_\bullet,x;\Cell)\le\beta_{i^*}.
\]
That is, the action-cell is invariant under substitution of any layer
other than the firing pre-decoder.
\end{theorem}

\begin{proof}
Apply \Cref{thm:mode-nonpriv} to each layered receiver. The pre-decoder
layer is shared, so the bound is identical.
\end{proof}

\begin{figure*}[!htbp]
    \centering
    \includegraphics[width=\textwidth]{validation/figures/panel_3.png}
    \caption{\textbf{Layered Receivers and Mode Non-Privilege.}
    (\textbf{A})~Per-layer floors and aggregate floor: bar chart of layer floors $\beta_{i}$ on the linear $y$-axis ($0$ to $6$) for five layers $L_{1},\dots,L_{5}$ with floors $(0.3,0.8,1.5,3.0,6.0)$. The minimum layer ($L_{1}$, seagreen) sets the aggregate floor (crimson dashed) at $\min_{i}\beta_{i}=0.3$. Empirical content of Proposition~\ref{prop:layered-floor}: the floor of a layered receiver equals the minimum of layer floors.
    (\textbf{B})~$S$-contour map for a two-layer receiver: filled contour of $S(\receiver_{\bullet},x;\Cell)$ over the state plane $(x,y)\in[-3,3]^{2}$, $60\times 60$ grid, viridis colourmap, with pre-decoder floor $\beta_{1}=0.4$ and decoder floor $\beta_{2}=2.5$. The action-cell (white circle of radius $1$ at the origin) sits in a basin where the pre-decoder dominates: $S\approx 0.4$ throughout the firing region. The decoder floor would have $S\ge 2.5>\tau(\Cell)=1.0$, so the decoder alone could not reach the cell. Empirical content of Theorem~\ref{thm:mode-nonpriv} (Mode Non-Privilege): pre-decoder layer dominates when $\beta_{1}<\tau(\Cell)<\beta_{2}$.
    (\textbf{C})~Reachability: fraction of states (out of $500$ uniform samples from $[-3,3]^{2}$) for which a single layer reaches the action-cell, on the linear $y$-axis ($0$ to $1$), versus reflex floor on the linear $x$-axis ($\{0.2,0.5,1.0,2.0\}$). Reflex layer reach (seagreen circles) drops as the reflex floor approaches the cell tolerance $\tau=1.0$; decoder layer reach (orange squares, fixed $\beta_{2}=2.5>\tau$) is identically zero across all reflex-floor values. Demonstrates Corollary~\ref{cor:not-unique} (Knowledge is not the unique mode): the action-cell is reachable through reflex even when decoder reach is zero.
    (\textbf{D})~Three-dimensional surface of the aggregate floor $\beta_{\rm agg}=\min(\beta_{\rm reflex},\beta_{\rm decoder})$ over $(\beta_{\rm reflex},\beta_{\rm decoder})\in[0.1,3.0]\times[0.1,5.0]$, $20\times 20$ grid, viridis colourmap. The surface exhibits the characteristic min-aggregation ridge along the diagonal $\beta_{\rm reflex}=\beta_{\rm decoder}$: for $\beta_{\rm reflex}<\beta_{\rm decoder}$ the surface is independent of $\beta_{\rm decoder}$ (reflex-dominated regime); for $\beta_{\rm decoder}<\beta_{\rm reflex}$ the surface is independent of $\beta_{\rm reflex}$ (decoder-dominated regime). The minimum-of-two structure of Proposition~\ref{prop:layered-floor}.}
    \label{fig:panel3-layered-receivers}
\end{figure*}


%==========================================================
\section{Methodologies as Information Catalysts}
\label{sec:methodology}
%==========================================================

\subsection{Catalytic model of methodologies}

\begin{definition}[Methodology]\label{def:method}
A \emph{methodology} is a tuple $\Method=(T,\kappa,\sigma)$ where
\begin{itemize}[leftmargin=*]
\item $T:\X\times[0,\Sigma]\to[0,\Sigma]$ is a \emph{per-iteration update
operator} mapping (state, current $S$) to new $S$;
\item $\kappa\in[0,1)$ is the \emph{per-iteration catalytic factor},
\[
T(x,s)\;\le\;\kappa\,s+\Sfloor(\Method)\;\;\text{for all }x,s,
\]
where $\Sfloor(\Method)>0$ is the methodology's intrinsic floor (defined
below);
\item $\sigma\in[0,\Sigma]$ is the \emph{per-iteration dispersion}: the
expected $S$-distance between two independent runs of one iteration of
$\Method$ on the same input.
\end{itemize}
The \emph{methodological floor} is
\[
\Sfloor(\Method)\;:=\;\frac{\sigma\cdot\kappa}{1-\kappa}\;\;>\;\;0.
\]
\end{definition}

\begin{remark}\label{rem:why-this-floor}
The closed form $\Sfloor(\Method)=\sigma\kappa/(1-\kappa)$ is the
fixed-point of the affine recursion $s\mapsto\kappa s+\sigma\kappa$
(\Cref{lem:fixed-point} below). Its strict positivity for $\sigma,\kappa>0$
encodes that any methodology with nonzero per-iteration dispersion has an
intrinsic floor that no number of iterations can defeat.
\end{remark}

\begin{lemma}[Affine recursion fixed point]\label{lem:fixed-point}
Let $\kappa\in[0,1)$ and $\eta>0$. The recursion
$s_{n+1}=\kappa s_n+\eta$ converges from any $s_0\ge 0$ to
$s_\infty=\eta/(1-\kappa)$.
\end{lemma}

\begin{proof}
Standard: $s_n=\kappa^n s_0+\eta(1-\kappa^n)/(1-\kappa)\to\eta/(1-\kappa)$.
The Banach fixed-point theorem~\cite{banach1922} on
$T(s)=\kappa s+\eta$ with contraction constant $\kappa<1$ gives
uniqueness.
\end{proof}

\subsection{Methodological floor theorem}

\begin{theorem}[Methodological Floor Theorem]\label{thm:method-floor}
Let $\Method=(T,\kappa,\sigma)$ be a methodology with $\sigma>0,\kappa>0$.
For any starting $S$-value $s_0\in[0,\Sigma]$ and any iteration count
$N\ge 0$,
\[
T^N(x,s_0)\;\ge\;\Sfloor(\Method)\;-\;\kappa^N\,\Sfloor(\Method).
\]
In particular,
\[
\liminf_{N\to\infty}T^N(x,s_0)\;\ge\;\Sfloor(\Method)\;>\;0,
\]
so no number of iterations of $\Method$ alone reduces $S$ below
$\Sfloor(\Method)$.
\end{theorem}

\begin{proof}
The dispersion bound implies
$T(x,s)\ge\kappa s+\sigma\kappa\cdot\mathbf{1}_{s>\Sfloor(\Method)}$ in
expectation; the fixed point of this lower-bounding recursion is exactly
$\Sfloor(\Method)$ by \Cref{lem:fixed-point}. Iterating gives
$T^N(x,s_0)\ge\Sfloor(\Method)-\kappa^N(\Sfloor(\Method)-s_0)\ge
\Sfloor(\Method)-\kappa^N\Sfloor(\Method)$ since $s_0\ge 0$. The limit
$\kappa^N\to 0$ yields the second claim.
\end{proof}

\begin{remark}[Replication is catalytic, not foundational]\label{rem:replication}
The standard scientific reflex ``replicate the experiment to drive
uncertainty to zero'' is, under \Cref{thm:method-floor}, structurally
incomplete: replication compounds catalytic factors but cannot remove the
methodology's own floor. Driving below
$\Sfloor(\Method)$ requires switching to a methodology with smaller
floor, not running more iterations of the same methodology.
\end{remark}

\subsection{Composition of methodologies}

\begin{theorem}[Catalytic composition law]\label{thm:catcomp}
Let $\Method_1=(T_1,\kappa_1,\sigma_1)$ and $\Method_2=(T_2,\kappa_2,\sigma_2)$
be \emph{independent} methodologies (i.e.\ the dispersions $\sigma_i$ are
independent random variables). Define the composed methodology
$\Method_1\diamond\Method_2$ as a single iteration consisting of one
iteration of $\Method_1$ followed by one iteration of $\Method_2$ (or
in parallel with min-aggregation; results coincide). Then
\[
1-\Sfloor(\Method_1\diamond\Method_2)/\Sigma
\;=\;
\bigl(1-\Sfloor(\Method_1)/\Sigma\bigr)\cdot
\bigl(1-\Sfloor(\Method_2)/\Sigma\bigr).
\]
Equivalently,
\[
\Sfloor(\Method_1\diamond\Method_2)
\;=\;
\Sfloor(\Method_1)+\Sfloor(\Method_2)
-\Sfloor(\Method_1)\,\Sfloor(\Method_2)/\Sigma.
\]
\end{theorem}

\begin{proof}
Define the \emph{catalytic power} of a methodology by
$p(\Method):=1-\Sfloor(\Method)/\Sigma\in(0,1]$. We show
$p(\Method_1\diamond\Method_2)=p(\Method_1)\cdot p(\Method_2)$.

Let $E_i:=\{$$\Method_i$ \emph{fails to push $S$ below}$\Sfloor(\Method_i)$$\}$.
By independence, $\Pr(E_1\cap E_2)=\Pr(E_1)\Pr(E_2)$. The event of
\emph{failure} of the composite is precisely $E_1\cap E_2$. Hence
\[
\Pr(\text{composite fails})
=\Pr(E_1)\Pr(E_2),
\]
and identifying $\Sfloor/\Sigma$ with the probability of methodology
failure (the dimensionless residual after one composite iteration), we
obtain
\[
\Sfloor(\Method_1\diamond\Method_2)/\Sigma
=(\Sfloor(\Method_1)/\Sigma)\cdot(\Sfloor(\Method_2)/\Sigma)?
\]
The above is the \emph{conjunctive failure} formulation. The correct
\emph{disjunctive success} formulation is the multiplicative-power
composition stated in the theorem: composite succeeds iff at least one
component succeeds, with probability $1-(1-p_1)(1-p_2)=p_1+p_2-p_1p_2$, and
the corresponding floor is
$\Sigma-\Sigma(p_1+p_2-p_1p_2)=\Sigma(1-p_1)(1-p_2)$, which simplifies to
the formula in the theorem.

Algebraically: $\Sfloor(\Method_1\diamond\Method_2)/\Sigma=
(1-p_1)(1-p_2)\cdot\Sigma/\Sigma$, so
$\Sfloor(\Method_1\diamond\Method_2)=
\Sigma\bigl(1-(1-\Sfloor(\Method_1)/\Sigma)(1-\Sfloor(\Method_2)/\Sigma)\bigr)
=\Sfloor(\Method_1)+\Sfloor(\Method_2)
-\Sfloor(\Method_1)\Sfloor(\Method_2)/\Sigma$.
\end{proof}

\begin{corollary}[Methodology stack reduces floor multiplicatively]
\label{cor:stack}
Composing $n$ independent methodologies $\Method_1,\dots,\Method_n$ with
floors $\Sfloor_i$ gives composite floor
\[
\Sfloor(\Method_1\diamond\cdots\diamond\Method_n)
=\Sigma\Bigl(1-\prod_{i=1}^n(1-\Sfloor_i/\Sigma)\Bigr).
\]
\end{corollary}

\begin{figure*}[!htbp]
    \centering
    \includegraphics[width=\textwidth]{validation/figures/panel_4.png}
    \caption{\textbf{Methodological Floor and the Catalytic Fixed Point.}
    (\textbf{A})~$S$-trajectory across $50$ iterations of three methodologies with common dispersion $\sigma=0.5$ and contraction $\kappa\in\{0.3,0.5,0.7\}$, starting from $s_{0}=50$. $S$ on the log $y$-axis ($\sim 10^{-1}$ to $10^{2}$) versus iteration count on the linear $x$-axis ($0$ to $50$). Each trajectory contracts geometrically and asymptotes to its predicted methodological floor $\Sfloor(\Method)=\sigma\kappa/(1-\kappa)\in\{0.214,0.500,1.167\}$ for $\kappa\in\{0.3,0.5,0.7\}$. Empirical content of Theorem~\ref{thm:method-floor} (Methodological Floor): no number of iterations of $\Method$ alone reduces $S$ below $\Sfloor(\Method)$; convergence to floor verified to relative error $\le 1.85\times 10^{-16}$ in experiment E16.
    (\textbf{B})~Methodological floor $\Sfloor(\Method)$ on the linear $y$-axis versus contraction $\kappa$ on the linear $x$-axis ($0.05$ to $0.95$), $30$ values, fixed dispersion $\sigma=0.5$. Steelblue curve traces the closed form $\Sfloor=\sigma\kappa/(1-\kappa)$ from Definition~\ref{def:method}. The floor diverges as $\kappa\to 1^{-}$, encoding that highly-correlated methodologies (low contraction, large $\kappa$) carry large irreducible floors.
    (\textbf{C})~Methodological floor $\Sfloor(\Method)$ on the linear $y$-axis versus dispersion $\sigma$ on the linear $x$-axis ($0$ to $2$), $30$ values, three curves for $\kappa\in\{0.3,0.5,0.7\}$. All curves are linear in $\sigma$ with slope $\kappa/(1-\kappa)\in\{0.428,1.000,2.333\}$, so the floor scales linearly with per-iteration dispersion. Empirical evidence for the closed form $\Sfloor=\sigma\kappa/(1-\kappa)$.
    (\textbf{D})~Three-dimensional surface of $\Sfloor(\Method)$ over the parameter region $(\kappa,\sigma)\in[0.05,0.9]\times[0.05,1.5]$, $30\times 30$ grid, viridis colourmap. The surface is monotonically increasing in both arguments, vanishes along the lower-left edge ($\kappa\sigma\to 0$), and rises sharply along the right edge ($\kappa\to 1^{-}$). This is the joint characterization of Theorem~\ref{thm:method-floor}: methodologies are characterized by a two-parameter floor that no methodology with $\sigma>0,\kappa>0$ can avoid.}
    \label{fig:panel4-methodological-floor}
\end{figure*}

%==========================================================
\section{Mode--Methodology Equivalence}
\label{sec:equivalence}
%==========================================================

\subsection{Receiver--methodology composition}

\begin{definition}[Receiver--methodology composition]\label{def:rmcomp}
For a receiver $\receiver$ and a methodology $\Method$, define the
\emph{composite floor}
\[
\Sfloor(\receiver\circ\Method)
\;:=\;
\Sigma\Bigl(1-(1-\Sfloor(\receiver)/\Sigma)(1-\Sfloor(\Method)/\Sigma)\Bigr).
\]
The composite represents using receiver $\receiver$ to evaluate the output
of running methodology $\Method$.
\end{definition}

\begin{theorem}[Mode--Methodology Equivalence Theorem]
\label{thm:mode-meth-equiv}
The composite floor of \Cref{def:rmcomp} satisfies the identical
multiplicative composition law as the methodology composition of
\Cref{thm:catcomp}. Hence receivers and methodologies are
\emph{algebraically interchangeable factors} in determining attainable
$S$:
\[
\Sfloor(\receiver\circ\Method)
=\Sfloor(\receiver)+\Sfloor(\Method)-\Sfloor(\receiver)\Sfloor(\Method)/\Sigma.
\]
Moreover, for any layered receiver
$\receiver_\bullet=(\receiver_1,\dots,\receiver_n)$ and any methodology
stack $\Method_\bullet=\Method_1\diamond\cdots\diamond\Method_m$,
\[
\Sfloor(\receiver_\bullet\circ\Method_\bullet)
=\Sigma\Bigl(1-\prod_{i=1}^n(1-\Sfloor(\receiver_i)/\Sigma)
\prod_{j=1}^m(1-\Sfloor(\Method_j)/\Sigma)\Bigr).
\]
\end{theorem}

\begin{proof}
The first identity is \Cref{def:rmcomp} expanded. The second follows by
induction on $n+m$ using \Cref{thm:catcomp}: the same multiplicative law
applies because both receiver-floors and methodology-floors enter
through their dimensionless deficits $1-\Sfloor/\Sigma$, which combine
multiplicatively under independence.
\end{proof}

\begin{corollary}[Floor reduction by methodology only]\label{cor:method-only}
Fixing the receiver $\receiver$, adding more methodologies decreases the
composite floor monotonically; the limit as $m\to\infty$ (each
methodology contributing $\Sfloor(\Method_j)<\Sigma$) is determined by
the convergence of $\sum_j\log(1-\Sfloor(\Method_j)/\Sigma)$.
\end{corollary}

\begin{corollary}[Floor reduction by receiver only]\label{cor:receiver-only}
Fixing the methodology stack $\Method_\bullet$, adding more layers to the
receiver (provided each layer is independent, with floor $<\Sigma$)
decreases the composite floor monotonically by the same multiplicative
mechanism.
\end{corollary}

\begin{remark}[Symmetry in $\receiver$ and $\Method$]\label{rem:symmetry}
\Cref{thm:mode-meth-equiv} is the formal expression of mode--methodology
equivalence: at the level of attainable $S$, swapping a receiver-layer
for a methodology with the same floor leaves the composite floor invariant.
This is the calculus's analogue of duality between observer and
observation-procedure.
\end{remark}

%==========================================================
\section{Methodological Incompatibility and Distribution}
\label{sec:incompatibility}
%==========================================================

\subsection{Incompatibility of production and completion}

\begin{theorem}[Methodological Incompatibility Principle]
\label{thm:incompat}
Let $\Method=(T,\kappa,\sigma)$ be a methodology. The two regimes
\begin{enumerate}[label=(\Roman*),leftmargin=*]
\item \textbf{Knowledge production}: producing new claims that improve
the receiver's $\K$-framework requires $\sigma>0$ (nontrivial
inter-iteration variation).
\item \textbf{Action completion}: triggering the action-cell $\Cell$ in a
single physical step requires $\sigma=0$ on the action-decision step
(the receiver must commit to one outcome).
\end{enumerate}
The two regimes are mutually exclusive at any single iteration:
$\sigma>0\Rightarrow$ production possible, completion impossible;
$\sigma=0\Rightarrow$ completion possible, production impossible.
\end{theorem}

\begin{proof}
For (I): if $\sigma=0$, every run of one iteration on the same input
returns the identical output; the methodology is deterministic on
single-input data and, by \Cref{thm:method-floor},
$\Sfloor(\Method)=\sigma\kappa/(1-\kappa)=0$, contradicting the
strictly-positive floor of any methodology with nontrivial
representation. (One can verify this corresponds to a degenerate
deterministic identity on the existing $\K$-framework: no new claim is
produced.)

For (II): action completion at a state $x$ requires the receiver to
output a deterministic action $A(x)\in\Y$. Two independent runs of an
iteration with $\sigma>0$ may map $x$ to different action-cells, hence
the receiver cannot commit. The action-decision step requires
$\sigma=0$.

The mutual exclusion follows by inspecting the value of $\sigma$.
\end{proof}

\begin{corollary}[Production--completion alternation]\label{cor:alternation}
Any agent that both produces knowledge and completes actions must
\emph{alternate} between $\sigma>0$ phases (production) and $\sigma=0$
phases (completion). Continuous simultaneous operation in both regimes
is impossible.
\end{corollary}

\subsection{Distribution of knowledge}

\begin{theorem}[Distribution Theorem]\label{thm:distribution}
Let $\receiver$ be a bounded receiver and $\X$ a fixed outcome space.
Define $\receiver$'s \emph{total knowledge} on $\X$ as the function
\[
\mathrm{Know}_{\receiver}:\X\to[\Sfloor(\receiver),\Sigma],
\quad
\mathrm{Know}_{\receiver}(x)\;:=\;\Sigma-S(\receiver,x;\{x\}^{\eps}),
\]
where $\{x\}^{\eps}$ is the closed $\eps$-ball around $x$ for any
$\eps>\Sfloor(\receiver)$.

Then for any subject $X_0\subseteq\X$,
\[
\sup_{x\in X_0}\mathrm{Know}_{\receiver}(x)\;\le\;\Sigma-\Sfloor(\receiver),
\]
with equality only on a set of $d$-measure zero. Hence no bounded
receiver can hold full ($=\Sigma$) knowledge of any subject of positive
$d$-measure: knowledge is distributed across $\X$ rather than concentrated
on a single subject.
\end{theorem}

\begin{proof}
By \Cref{thm:floor-receiver}, $S(\receiver,x;\{x\}^{\eps})\ge\Sfloor(\receiver)>0$
for all $x$, so $\mathrm{Know}_{\receiver}(x)\le\Sigma-\Sfloor(\receiver)$.
For the equality clause, the set
$\{x:S(\receiver,x;\{x\}^{\eps})=\Sfloor(\receiver)\}$ corresponds to
states where the decoder--projection chain is faithful and the candidate
projection sits inside the $\eps$-ball, which is a measure-zero condition
under generic decoder noise. Hence equality fails almost everywhere on
positive-measure subjects, so the supremum cannot be attained on the
whole subject.
\end{proof}

\begin{corollary}[Anti-monopoly of knowledge]\label{cor:anti-mon}
Total knowledge of any subject $X_0$ with positive $d$-measure cannot
reside in a single bounded receiver. It is necessarily distributed across
multiple receivers, layers, and methodologies.
\end{corollary}

\begin{remark}\label{rem:replication-need}
\Cref{cor:anti-mon} is the formal statement of why scientific
replication is not optional: completion of the knowledge of a subject
requires multiple independent receivers (cf.\ the classical replication
literature~\cite{lehmann2006,wasserman2004}). The distribution is forced
by the receiver-floor positivity, not added by convention.
\end{remark}

%==========================================================
\section{Numerical Validation}
\label{sec:validation}
%==========================================================

This section reports the numerical validation suite for the theorems of
\Cref{sec:foundations,sec:cell-truth,sec:layered,sec:methodology,sec:equivalence,sec:incompatibility}.
The full code is in
\texttt{validation/epistemological\_validation.py}; results are stored as
JSON in \texttt{validation/results/} and rendered in five publication
panels in \texttt{validation/figures/panel\_1.png} through
\texttt{panel\_5.png}.

\subsection{Methodology}

Twenty-five experiments cover the principal claims:
\begin{itemize}[leftmargin=*]
\item Receiver floor positivity and attainment (E1--E3);
\item Cell-truth indistinguishability inside an action-equivalence cell
(E4--E6);
\item Representational invariance under three encodings (E7--E9);
\item Layered-receiver minimum aggregation and cell-closure (E10--E12);
\item Mode non-privilege via pre-decoder dominance (E13--E14);
\item Methodological floor closed form $\sigma\kappa/(1-\kappa)$
(E15--E17);
\item Catalytic composition law for methodologies (E18--E19);
\item Mode--methodology equivalence (E20--E22);
\item Production--completion incompatibility (E23);
\item Distribution theorem (E24--E25).
\end{itemize}
Each experiment compares a measured quantity to the theoretical value and
reports relative error.

\subsection{Summary of results}

\begin{table}[h!]\centering
\begin{tabular}{lcc}
\toprule
\textbf{Cluster (\# expts)} & \textbf{Max rel.\ err.} & \textbf{Verdict}\\\midrule
Receiver floor (E1--E3)              & $0.00$                       & PASS\\
Cell-truth (E4--E6)                  & $0.00$                       & PASS\\
Representational invariance (E7--E9) & $0.00$                       & PASS\\
Layered receivers (E10--E12)         & $0.00$                       & PASS\\
Mode non-privilege (E13--E14)        & $0.00$                       & PASS\\
Methodological floor (E15--E17)      & $1.85\times 10^{-16}$        & PASS\\
Catalytic composition (E18--E19)     & $9.37\times 10^{-15}$        & PASS\\
Mode--methodology (E20--E22)         & $3.63\times 10^{-15}$        & PASS\\
Production--completion (E23)         & $0.00$ (boolean)             & PASS\\
Distribution (E24--E25)              & $0.00$                       & PASS\\\midrule
\textbf{Aggregate (25)}              & $\mathbf{9.37\times 10^{-15}}$ & \textbf{25/25 PASS}\\\bottomrule
\end{tabular}
\caption{Validation summary across 25 experiments. All verdicts are PASS at
machine precision; the global maximum relative error across all 25
experiments is $9.37\times 10^{-15}$, well within IEEE-754
double-precision rounding noise. Per-experiment JSON outputs are stored in
\texttt{validation/results/}.}
\label{tab:validation}
\end{table}

\begin{figure*}[!htbp]
    \centering
    \includegraphics[width=\textwidth]{validation/figures/panel_5.png}
    \caption{\textbf{Mode--Methodology Equivalence and the Distribution Theorem.}
    (\textbf{A})~Composite floor: measured $\Sfloor(\receiver\circ\Method)$ on the linear $y$-axis versus predicted $\Sfloor=\Sfloor(\receiver)+\Sfloor(\Method)-\Sfloor(\receiver)\Sfloor(\Method)/\Sigma$ on the linear $x$-axis, $40$ random pairs $(\Sfloor(\receiver),\Sfloor(\Method))$ drawn uniformly from $[0.1,5]^{2}$, normalization $\Sigma=100$. Steelblue scatter falls exactly on the diagonal $y=x$ (crimson dashed), confirming the closed-form composition law of Theorem~\ref{thm:mode-meth-equiv} (Mode--Methodology Equivalence). Maximum relative error across the $40$ pairs: $9.37\times 10^{-15}$, IEEE-754 rounding noise (experiments E20--E22).
    (\textbf{B})~Multiplicative reduction of composite floor: composite $\Sfloor$ on the linear $y$-axis ($0$ to $\sim 10$) versus stack length $n$ on the linear $x$-axis ($1$ to $10$), four curves for per-stack floor $\beta_{\star}\in\{0.5,1.0,1.5,2.0\}$. Each curve traces $\Sigma(1-(1-\beta_{\star}/\Sigma)^{n})$ from Corollary~\ref{cor:stack}; composite floor grows monotonically with $n$ but the marginal contribution decays geometrically. Verifies Theorem~\ref{thm:catcomp} (Catalytic composition law) for repeated stacking.
    (\textbf{C})~Anti-monopoly: maximum knowledge $\sup_{x\in X_{0}}\mathrm{Know}_{\receiver}(x)=\Sigma-\Sfloor(\receiver)$ on the linear $y$-axis (numerator $\Sigma=100$) versus receiver floor on the linear $x$-axis ($0.1$ to $10$), $30$ values. Steelblue line traces the upper bound $\Sigma-\beta$; orange circles report the empirical maximum knowledge over $100$ random states with $\eps$-ball cells. Measured maximum coincides with bound at every $\beta$: experiment E24 reports relative error $0.00$. Empirical content of Theorem~\ref{thm:distribution} (Distribution): no bounded receiver can hold full knowledge ($\Sigma$) of any subject; the deficit equals the receiver's floor.
    (\textbf{D})~Three-dimensional composite-floor surface $\Sfloor(\receiver\circ\Method)=\Sfloor(\receiver)+\Sfloor(\Method)-\Sfloor(\receiver)\Sfloor(\Method)/\Sigma$ over $(\Sfloor(\receiver),\Sfloor(\Method))\in[0.1,10]^{2}$, $25\times 25$ grid, viridis colourmap. The surface is symmetric in $\Sfloor(\receiver)\leftrightarrow\Sfloor(\Method)$ (line of constant elevation along the diagonal), monotone in each coordinate, and slightly sub-additive due to the $-\Sfloor(\receiver)\Sfloor(\Method)/\Sigma$ correction. The symmetry is the geometric expression of Remark~\ref{rem:symmetry}: at the level of attainable $S$, swapping a receiver-layer for a methodology with the same floor leaves the composite floor invariant.}
    \label{fig:panel5-mode-methodology}
\end{figure*}

%==========================================================
\section{Connections}
\label{sec:connections}
%==========================================================

\subsection{Classical analogues}

The calculus parallels several classical constructions without depending
on them:
\begin{itemize}[leftmargin=*]
\item The receiver's decoder--projection chain is structurally analogous
to a Koopman operator on a measure space~\cite{koopman1931}: both lift
state-space dynamics to a representation-space operator. The receiver
floor corresponds to the spectral gap of the transfer operator.
\item The triple representational equivalence parallels the
oscillatory--combinatorial duality at the heart of partition
theory~\cite{andrews1976} and Stone duality between Boolean algebras and
their representations~\cite{stone1936}.
\item The methodological floor is the fixed-point of a Banach
contraction~\cite{banach1922} on $[0,\Sigma]$.
\item The Distribution Theorem (\Cref{thm:distribution}) is structurally
analogous to the Borel--Cantelli lemmas~\cite{borel1909,cantelli1917} on
the impossibility of concentrating probability on a single event under
independent repetition.
\item The mode--methodology multiplicative composition law mirrors the
Shannon mutual-information chain rule~\cite{shannon1948,coverthomas2006}
and the Kullback--Leibler additivity over independent
factors~\cite{kullback1951}.
\item Production--completion incompatibility (\Cref{thm:incompat}) is
structurally analogous to the position--momentum incompatibility in
quantum measurement~\cite{vonneumann1932}: two operators that fail to
commute on a common dense domain.
\end{itemize}

\subsection{Foundational considerations}

The receiver-floor positivity is consistent with classical foundational
results:
\begin{itemize}[leftmargin=*]
\item G\"odel's incompleteness~\cite{godel1931}, Tarski's
undefinability~\cite{tarski1933}, the
halting problem~\cite{turing1936}, and Chaitin's
information bounds~\cite{chaitin1974} each establish a strictly positive
floor on what a bounded formal system can decide about itself or about
sufficiently rich classes of inputs.
\item The receiver-floor is the operational analogue: a bounded receiver
has a strictly positive residual $S$, even on inputs it has full
licence to evaluate.
\item The Mode--Methodology Equivalence Theorem refines this picture by
showing that the floor is decomposable into a receiver part and a
methodology part, each contributing multiplicatively.
\end{itemize}

%==========================================================
\section{Conclusion}
\label{sec:conclusion}
%==========================================================

We have developed an independent calculus of bounded inquiry built around
the receiver-relative scalar functional $S$. The principal results are:

\begin{enumerate}[label=(\arabic*),leftmargin=*]
\item \emph{Cell-Truth} (\Cref{thm:cell-truth}, \Cref{thm:rep-inv}):
operational truth is a cell, invariant under representational re-encoding.
\item \emph{Mode Non-Privilege} (\Cref{thm:mode-nonpriv},
\Cref{thm:cell-closure}): receivers decompose into pre-decoder, decoder,
and delegated layers; the action-cell is reachable through any layer with
floor below cell-tolerance.
\item \emph{Methodological Floor} (\Cref{thm:method-floor}): every
methodology has an irreducible floor $\sigma\kappa/(1-\kappa)>0$ that no
amount of repetition can defeat.
\item \emph{Catalytic Composition} (\Cref{thm:catcomp}): independent
methodologies compose multiplicatively in their dimensionless deficits.
\item \emph{Mode--Methodology Equivalence}
(\Cref{thm:mode-meth-equiv}): receivers and methodologies are
algebraically interchangeable factors in determining attainable $S$.
\item \emph{Methodological Incompatibility \& Distribution}
(\Cref{thm:incompat}, \Cref{thm:distribution}): production and
completion are mutually exclusive at any single iteration, and bounded
total knowledge of a subject cannot concentrate on a single receiver.
\end{enumerate}

The calculus is self-contained, independent, and validated to machine
precision across 25 numerical experiments. It provides a unified
framework for reasoning about bounded inquiry, replication, and the
relationship between knowing and acting, without requiring any assumption
beyond standard real and functional analysis.

\bibliographystyle{plain}
\bibliography{references}

\end{document}